\begin{musica}{Moçoila}{Jorge Prendas, Mário Alves}{Tuna Académica da Universidade Portucalense}
Olá, \ac{C}{lin}da, mo\ac{C/B}{çoi}la,
Que \ac{Am}{tra}zes p'ra \ac{G}{mim}.
P'ra além \ac{F}{da} maldi\ac{C}{ção},
De seres \ac{F}{tão} bela as\ac{G}{sim}.
Trazes \ac{F}{bei}jos e \ac{Em}{le}vas-me, o \ac{Dm}{tro}co no \ac{G}{fim}. \ac{F G C}{}_ _ _

Tapa os lábios com seda,
Não vá ver-te alguém.
E ficar encantado,
Com eles, meu bem.
E se a seda chegar,
Tapa os olhos também.

\ac{Am}{O}lha que o \ac{Em}{dia} \ac{F}{não} é se\ac{C}{gu}ro,
\ac{Am}{Guar}da-te em \ac{Em}{ca}sa.
\ac{F}{Sem}pre que o sol nas\ac{G}{cer},
\ac{Am}{O}lha que a \ac{Em}{noi}te \ac{F}{tem} mau fei\ac{C}{tio}.
\ac{Am}{Não} te aven\ac{Em}{tu}res \ac{F}{sem}pre que anoite\ac{G}{cer}.

Irei \ac{C}{ver}-te à noi\ac{C/B}{ti}nha,
P'ra o \ac{A#dim7}{po}vo não \ac{A7}{ver}.
Porque as\ac{F}{sim} pelo es\ac{C}{cu}ro,
Nem \ac{Am}{som}bra vou \ac{G}{ter}.
Ele há \ac{F}{som}bras p'ra \ac{Em}{tu}do,
Nem \ac{Am}{quei}ras sa\ac{G}{ber}.
Olá, \ac{F}{lin}da, mo\ac{Am}{çoi}la,
Pre\ac{F}{pa}ra-te, \ac{G}{vou} ba\ac{C}{ter}!

\posicoes{C: x32010, C/B: x22010, Dm: xx0231, Em: 022000, F: 133211, G: 320033,
  Am: x02210, A7: x02020, A#dim7: x12020}
\end{musica}
